\section{System General Perspective}

\subsection{Scope and Objectives}
According to Project Proposal 13, the SWAID project is an interactive standing wave demonstrator using sound. The specific objectives that impact the PR subsystem include:
\begin{itemize}
    \item Modulation of generated images through background music and hand-tracking information.
    \item Synchronized lighting to increase the demonstrator's attractiveness.
\end{itemize}

\subsection{Global System Vision}
SWAID is composed of three articulated subsystems: CP, PR, and HI. The PR subsystem is responsible for the physical excitation of the plate and managing the resonance required for the formation of Chladni patterns, receiving control signals from the interface subsystem.

\subsection{Features}
\begin{itemize}
    \item Controlled excitation of the plate via acoustic transducers.
    \item Dynamic frequency adjustment to obtain different modal patterns.
    \item Response to tracking inputs (coming from the HI subsystem).
\end{itemize}

\subsection{User Characteristics}
Students and the general public at scientific outreach events. The system must be visually appealing and easy to operate.

\subsection{Constraints}
\begin{itemize}
    \item Limited budget for hardware components.
    \item Portability.
\end{itemize}